\subsection{Contactos}

Para cualquier consulta o problema relacionado con el servicio cubierto por este SLA, se establecen los siguientes contactos de soporte organizados por niveles de atención:

\subsubsection*{Nivel 1 \\ Mesa de Ayuda de Aulas}
\begin{itemize}
	\item \textbf{Nombre:} Ing. David Mata Guerra\\
    	\textbf{Puesto:} Analista de Soporte de Primer Nivel (Mesa de Ayuda)\\
    	\textbf{Correo:} \href{mailto:mesa.ayuda.aulas@universidad.edu}{mesa.ayuda.aulas@universidad.edu}\\
    	\textbf{Teléfono:} +52 (442) 192 1200 Ext. 3101
\end{itemize}

Responsable de recibir y registrar todos los reportes de fallas y solicitudes de servicio, realizar el diagnóstico inicial y resolver incidentes básicos relacionados con conectividad de usuario, autenticación y acceso a recursos.

\subsubsection*{Nivel 2 \\ Soporte Especializado de Redes}
\begin{itemize}
	\item \textbf{Nombre:} Mtro. Gabriel Juárez Ramírez\\
    	\textbf{Puesto:} Ingeniero de Redes Campus Aulas\\
    	\textbf{Correo:} \href{mailto:soporte.redes.aulas@universidad.edu}{soporte.redes.aulas@universidad.edu}\\
    	\textbf{Teléfono:} +52 (442) 192 1200 Ext. 3205
\end{itemize}

Encargado de atender incidentes escalados desde el Nivel 1 que involucren equipo de conmutación, enlaces troncales, VLANs, puntos de acceso inalámbricos o configuración del firewall perimetral.

\subsubsection*{Nivel 3 \\ Coordinación de Infraestructura y Escalamiento a Proveedores}
\begin{itemize}
	\item \textbf{Nombre:} Mtro. Alejandro Macías Fonseca\\
    	\textbf{Puesto:} Coordinador de Infraestructura de Redes\\
    	\textbf{Correo:} \href{mailto:infraestructura.redes@universidad.edu}{infraestructura.redes@universidad.edu}\\
    	\textbf{Teléfono:} +52 (442) 192 1200 Ext. 3302
\end{itemize}

Responsable de la gestión de incidentes críticos, decisiones de cambio mayor, coordinación con el proveedor de servicios de Internet y con fabricantes (por ejemplo, soporte de Fortinet o Aruba) cuando se requiera soporte de tercer nivel.\\[0.5em]

Como contacto alterno de nivel ejecutivo se designa a \textbf{Alexis Emilio Cárdenas Camacho}, responsable de supervisar el cumplimiento general del SLA y aprobar acciones extraordinarias cuando sea necesario.

\subsubsection*{Canales Generales de Atención}
Además de los contactos anteriores, se disponen los siguientes canales generales para el registro y seguimiento de casos:
\begin{itemize}
    \item \textbf{Portal de Soporte en Línea:} \url{https://soporte.universidad.edu/redes}
    \item \textbf{Teléfono de la Mesa de Ayuda:} +52 (442) 192 1200 Ext. 3100
\end{itemize}

El flujo de escalamiento se inicia siempre en el Nivel 1 y progresa a los niveles superiores únicamente cuando el incidente no pueda ser resuelto en el nivel previo o cuando, por su impacto, sea clasificado como incidente mayor según lo definido en este SLA.
